\documentclass[14pt]{extarticle}
\usepackage{preamble}
\geometry{letterpaper, landscape, margin=0.5in}
\usepackage{qrcode}
\renewcommand{\answer}[1]{}
\setlength{\parskip}{4pt}

\begin{document}
\pagestyle{empty}

\daybanner{Unit 2 \textperiodcentered\ Topic 2.10 \textperiodcentered\ TODAY'S PROBLEM}{Unusual Does Not Mean Impossible}

\noindent\begin{minipage}[t]{0.57\linewidth}
\vspace{0pt}
Suppose the same 70 percent free-throw shooter takes 100 attempts under comparable conditions and makes 58. Let $X$ be the number of makes if the career-rate model still applies.\par\smallskip \textbf{Coach:} ``Anything under 60 makes is impossible for a 70 percent shooter, so 58 proves the career rate no longer applies.''\par \textbf{Analyst:} ``A result can be unusual without being impossible; compare 58 with the model's mean and standard deviation.''\par The binomial model gives a chance of about $0.0072$ for 58 or fewer makes; use this provided value.\par
\end{minipage}\hfill
\begin{minipage}[t]{0.40\linewidth}
\vspace{0pt}
\begin{center}
\textbf{Career-rate model}\par\smallskip
\begin{tabular}{|l|c|c|c|}
\hline
\textbf{Set} & \textbf{Attempts} & \textbf{Rate} & \textbf{Makes} \\ \hline
\textbf{Tonight} & 100 & 0.70 & 58 \\ \hline
\end{tabular}
\end{center}
\end{minipage}

\begin{firsttakebox}{FIRST TAKE (5 min, on your sheet)}
Whose statement is more defensible, the coach's or the analyst's? Name one quantity you would check.
\end{firsttakebox}

\begin{description}[leftmargin=1.15in, labelwidth=1in, itemsep=2pt, topsep=2pt]
  \item[\dokbadge{1}~(a)] State $n$ and $p$ for this model.
  \item[\dokbadge{2}~(b)] Calculate the binomial mean and standard deviation. Interpret the mean for many sets of 100 shots.
  \item[\dokbadge{3}~(c)] Whose statement is supported? Compare 58 with the mean and standard deviation, and use the given probability to distinguish unusual from impossible. Explain what the coach's claim could make a reader wrongly believe. Write 3--4 sentences.
\end{description}

\vfill
\noindent\begin{minipage}{0.84\linewidth}
\textbf{Video + follow-along:} \href{https://robjohncolson.github.io/apstats-live-worksheet/u4_lesson10-12_live.html}{\texttt{\detokenize{u4_lesson10-12_live.html}}} (scan the code)\par
\textbf{Finish (a)--(c) after the video. Turn in your sheet.}
\end{minipage}\hfill
\qrcode[height=0.75in]{https://robjohncolson.github.io/apstats-live-worksheet/u4_lesson10-12_live.html}

\end{document}
